\section{Derivatives of elementary functions}

Many functions appear so often that their derivatives should become familiar. The goal is not only to memorize a table, but to recognize the local behavior of standard functions.

\subsection*{Exponential and logarithmic functions}

The exponential function \(e^x\) has a special role: it is its own derivative.

\begin{theorembox}
\[
\frac{d}{dx}e^x=e^x.
\]
For \(x>0\),
\[
\frac{d}{dx}\ln x=\frac1x.
\]
\end{theorembox}

\begin{examplebox}
Differentiate
\[
f(x)=3e^x-2\ln x.
\]
Using the rules,
\[
f'(x)=3e^x-\frac{2}{x},
\]
for \(x>0\).
\end{examplebox}

The logarithm already reminds us to watch the domain. The formula \((\ln x)'=1/x\) is used where \(x>0\).

\subsection*{Trigonometric functions}

The basic trigonometric derivatives are:
\[
\frac{d}{dx}\sin x=\cos x,
\]
\[
\frac{d}{dx}\cos x=-\sin x,
\]
\[
\frac{d}{dx}\tan x=\frac{1}{\cos^2 x}.
\]
The last expression is also written as \(\sec^2 x\), but this course will mostly use \(1/\cos^2 x\).

\begin{examplebox}
Differentiate
\[
f(x)=2\sin x-5\cos x.
\]
Then
\[
f'(x)=2\cos x+5\sin x.
\]
The sign changes because the derivative of \(\cos x\) is \(-\sin x\).
\end{examplebox}

\subsection*{Inverse trigonometric functions}

Some later examples require inverse trigonometric functions. The basic derivatives are:
\[
\frac{d}{dx}\arctan x=\frac{1}{1+x^2},
\]
\[
\frac{d}{dx}\arcsin x=\frac{1}{\sqrt{1-x^2}},
\qquad -1<x<1,
\]
\[
\frac{d}{dx}\arccos x=-\frac{1}{\sqrt{1-x^2}},
\qquad -1<x<1.
\]

\begin{examplebox}
Differentiate
\[
f(x)=x\arctan x.
\]
This expression is a product, so the product rule will be needed. Once that rule is introduced, we obtain
\[
f'(x)=\arctan x+\frac{x}{1+x^2}.
\]
\end{examplebox}

This example previews an important point: knowing elementary derivatives is not enough. We must also know how functions are combined.

\subsection*{A compact table}

\begin{center}
\begin{tabular}{ll}
\toprule
Function & Derivative \\
\midrule
\(c\) & \(0\) \\
\(x^p\) & \(p x^{p-1}\) \\
\(e^x\) & \(e^x\) \\
\(\ln x\) & \(1/x\) \\
\(\sin x\) & \(\cos x\) \\
\(\cos x\) & \(-\sin x\) \\
\(\tan x\) & \(1/\cos^2 x\) \\
\(\arctan x\) & \(1/(1+x^2)\) \\
\(\arcsin x\) & \(1/\sqrt{1-x^2}\) \\
\(\arccos x\) & \(-1/\sqrt{1-x^2}\) \\
\bottomrule
\end{tabular}
\end{center}

A derivative table is useful only if it is combined with careful reading of the expression. Before applying a formula, identify the function and its domain.

\begin{summarybox}
When using elementary derivatives, ask:
\begin{itemize}
\item Which elementary function appears?
\item Is the expression exactly elementary, or is it composed with another function?
\item What domain restrictions are present?
\item Are product, quotient, or chain rules also needed?
\end{itemize}
\end{summarybox}

Elementary derivatives are building blocks. Most real examples are built by combining them.